\chapter{Ciencia, salud y control social}
El surgimiento de nuevas universidades y la circulación de tratados médicos transformaron la comprensión del cuerpo y la enfermedad. Arrizabalaga estudia la llamada ``gran plaga francesa'' y muestra cómo la sífilis obligó a replantear las prácticas de higiene urbana y las medidas de cuarentena.\cite{Arrizabalaga1997} Estas políticas sanitarias conectan el saber médico con el incipiente aparato burocrático de las ciudades renacentistas.

\begin{figure}[H]
    \centering
    % Reemplazar por la imagen recortada correspondiente
    \includegraphics[width=0.7\textwidth]{FIGURES/tratado_medico_placeholder}
    \caption{Grabado médico sobre tratamientos de la sífilis ("Jon Arrizabalaga - The great pox: the French disease in Renaissance Europe (1997, Yale University Press) - libgen.li.pdf", p. XX).}
    \label{fig:tratado-medico}
\end{figure}

El control social también se vio atravesado por debates sobre la moral sexual y el rol del Estado. Saslow documenta cómo la figura de Ganímedes sirvió para discutir la homosexualidad dentro del arte y la literatura, revelando tensiones entre deseo, censura y patronazgo.\cite{Saslow1986} Estos debates se relacionan con la construcción de cuerpos normativos y con la forma en que la justicia urbana reguló el comportamiento público.
